% Resumo e abstract.
% ESTADO: rascunho. Os números citados já estão medidos, mas o resumo é a ÚLTIMA coisa a
% reescrever — ele precisa refletir o capítulo 4 fechado, inclusive usabilidade.

\chapter*{Resumo}
\addcontentsline{toc}{chapter}{RESUMO}

A interpretação de análises de solo e a derivação de recomendações de calagem e adubação
exigem conhecimento técnico e são sujeitas a erro, com consequências financeiras diretas
para o produtor. No Rio Grande do Sul e em Santa Catarina, os critérios estão
consolidados no Manual de Calagem e Adubação da CQFS-RS/SC, mas as ferramentas
computacionais disponíveis concentram-se na recomendação de correção, sem estimar a
aptidão da cultura escolhida frente às condições do solo da área; os sistemas que avaliam
aptidão, por sua vez, operam em escala regional e não se aplicam a uma análise
individual. Este trabalho apresenta o SIRAS, sistema especialista baseado em regras que
integra as duas perspectivas a partir de uma única análise de solo, cobrindo as 61
culturas e grupos de culturas para os quais o Manual publica tabelas completas de
adubação. O sistema interpreta os atributos informados, calcula a necessidade de calcário
pelo índice SMP e as doses de nitrogênio, fósforo e potássio, indica o adubo formulado
que atende à recomendação pelo procedimento do próprio Manual, e estima a aptidão edáfica
em dois cenários --- o estado atual do solo e o projetado após as correções recomendadas.
Toda saída é rastreável à regra, à tabela e à página que a produziram. A avaliação
combinou verificação automatizada, concordância com um oráculo calculado manualmente,
conformidade do módulo de aptidão frente a uma especificação aprovada antes da
implementação, e teste de usabilidade com usuários do público-alvo. Os casos com oráculo
apresentaram concordância integral, e a conformidade da aptidão foi de 100\% em classe e
em fator determinante. \pendente{acrescentar o resultado da usabilidade em uma frase}
Conclui-se que a automatização dos critérios é viável e que a integração entre
recomendação e aptidão, a partir de dados individuais de solo, representa contribuição
frente às ferramentas disponíveis para o contexto regional.

\textbf{Palavras-chave:} sistema especialista; análise de solo; calagem e adubação;
aptidão edáfica; sistemas baseados em regras.

\chapter*{Abstract}
\addcontentsline{toc}{chapter}{ABSTRACT}

Interpreting soil analyses and deriving liming and fertilization recommendations requires
technical expertise and is error-prone, with direct financial consequences for farmers.
In southern Brazil, the criteria are consolidated in the CQFS-RS/SC Liming and
Fertilization Manual, yet available software focuses on correction recommendations
without assessing the suitability of the chosen crop for the soil conditions of the
area; conversely, land suitability systems operate at a regional scale and do not apply
to an individual soil analysis. This work presents SIRAS, a rule-based expert system that
integrates both perspectives from a single soil analysis, covering the 61 crops and crop
groups for which the Manual publishes complete fertilization tables. The system
interprets the reported attributes, computes lime requirement from the SMP buffer index
along with nitrogen, phosphorus and potassium rates, indicates the compound fertilizer
that meets the recommendation following the Manual's own procedure, and estimates
\emph{edaphic} suitability under two scenarios --- the soil as it is, and as projected
after the recommended corrections. Every output is traceable to the rule, table and page
that produced it. Evaluation combined automated verification, agreement against a
manually computed oracle, conformance of the suitability module to a specification
approved before implementation, and a usability test with target users. Cases with an
oracle showed full agreement, and suitability conformance reached 100\% for both class
and limiting factor. \pendente{add the usability result in one sentence} We conclude that
automating these criteria is feasible, and that integrating recommendation and
suitability from individual soil data constitutes a contribution over the tools available
in the regional context.

\textbf{Keywords:} expert system; soil analysis; liming and fertilization; edaphic
suitability; rule-based systems.
